% !TEX root = ../main.tex
% ==============================================================================
% ĐỀ CƯƠNG ĐỒ ÁN TỐT NGHIỆP ĐẠI HỌC
% Chuẩn Khoa CNTT - Bộ môn Hệ thống Thông tin (HUMG)
% Theo mẫu Mau-De-cuong-DATN.docx & Hieu.pdf (Chuẩn đúng 5 trang: vi, vii, viii, ix, x)
% ==============================================================================

\thispagestyle{adminform}
\phantomsection
\addcontentsline{toc}{chapter}{ĐỀ CƯƠNG ĐỒ ÁN TỐT NGHIỆP ĐẠI HỌC}

{\setstretch{1.15}

% ------------------------------------------------------------------------------
% TRANG 1 / 5 (Trang vi): Tiêu ngữ, Bảng SV, Bảng GVHD, 1. Tên đề tài, 2. Lý do (1-2)
% ------------------------------------------------------------------------------
\noindent
\begin{tabular}{@{}p{6.8cm} >{\centering\arraybackslash}p{8.2cm}@{}}
  \centering \textbf{KHOA CÔNG NGHỆ THÔNG TIN} \newline \textbf{BỘ MÔN HỆ THỐNG THÔNG TIN} & 
  \textbf{CỘNG HÒA XÃ HỘI CHỦ NGHĨA VIỆT NAM} \newline
  \textbf{Độc lập -- Tự do -- Hạnh phúc} \newline
  \rule{3.5cm}{0.5pt}
\end{tabular}

\vspace{0.35cm}
\begin{center}
    {\fontsize{15pt}{19pt}\selectfont \textbf{ĐỀ CƯƠNG ĐỒ ÁN TỐT NGHIỆP ĐẠI HỌC}}
\end{center}
\vspace{0.2cm}

\noindent \textbf{Sinh viên thực hiện:}
\vspace{0.08cm}

\noindent
\begin{table}[h!]
\centering
\fontsize{9.5pt}{12pt}\selectfont
\setlength{\tabcolsep}{2.5pt}
\renewcommand{\arraystretch}{1.2}
\begin{tabular}{|>{\centering\arraybackslash}p{0.55cm}|>{\centering\arraybackslash}p{1.85cm}|>{\raggedright\arraybackslash}p{2.75cm}|>{\centering\arraybackslash}p{2.55cm}|>{\centering\arraybackslash}p{1.85cm}|>{\raggedright\arraybackslash}p{4.25cm}|}
\hline
\textbf{TT} & \textbf{Mã sinh viên} & \textbf{Họ và tên} & \textbf{Lớp} & \textbf{Điện thoại} & \textbf{E-mail} \\ \hline
1 & 2221050359 & Đàm Mạnh Dũng & DCCTCT67\_09B & 0978461865 & dammanhdungvn@gmail.com \\ \hline
\end{tabular}
\end{table}

\vspace{-0.25cm}
\noindent \textbf{Giảng viên hướng dẫn:}
\vspace{0.08cm}

\noindent
\begin{table}[h!]
\centering
\fontsize{9.5pt}{12pt}\selectfont
\setlength{\tabcolsep}{2.5pt}
\renewcommand{\arraystretch}{1.2}
\begin{tabular}{|>{\centering\arraybackslash}p{0.55cm}|>{\centering\arraybackslash}p{1.85cm}|>{\raggedright\arraybackslash}p{2.75cm}|>{\centering\arraybackslash}p{2.55cm}|>{\centering\arraybackslash}p{1.85cm}|>{\raggedright\arraybackslash}p{4.25cm}|}
\hline
\textbf{TT} & \textbf{Mã cán bộ} & \textbf{Họ và tên} & \textbf{Đơn vị} & \textbf{Điện thoại} & \textbf{E-mail} \\ \hline
1 & 0802-09 & TS. Vũ Lan Phương & Khoa CNTT & 0984223379 & vulanphuong@humg.edu.vn \\ \hline
\end{tabular}
\end{table}

\vspace{0.15cm}
\noindent \textbf{1. Tên đề tài}
\vspace{0.08cm}

Xây dựng nền tảng AI hỗ trợ nghiên cứu và viết bài báo khoa học VinUni dựa trên Agentic RAG, truy xuất tri thức học thuật và kiểm chứng trích dẫn.

\vspace{0.25cm}
\noindent \textbf{2. Lý do chọn đề tài}
\vspace{0.08cm}

Trong bối cảnh cuộc Cách mạng công nghiệp lần thứ tư, Trí tuệ nhân tạo (AI) và các Mô hình ngôn ngữ lớn (LLM) đang tạo ra bước chuyển dịch mang tính đột phá trong hoạt động nghiên cứu khoa học và giáo dục đại học. Khả năng tổng hợp thông tin, gợi ý lập luận và hỗ trợ viết bài của AI đã mở ra cơ hội to lớn giúp các nhà nghiên cứu rút ngắn đáng kể thời gian chuẩn bị công bố khoa học.

Đặc biệt, tại môi trường học thuật chất lượng cao như Trường Đại học VinUni, các giảng viên và học giả đối mặt với khối lượng lớn các ấn phẩm quốc tế đa nguồn cần đọc hiểu, tổng hợp và thẩm định. Việc tiếp cận nhanh chóng với các phương pháp tiên tiến, đồng thời duy trì năng suất công bố khoa học chuẩn IEEE/ACM là nhu cầu thường trực của các nhóm nghiên cứu.

\newpage
\thispagestyle{adminform}

% ------------------------------------------------------------------------------
% TRANG 2 / 5 (Trang vii): Lý do (tiếp), Giải pháp & Mục tiêu (Tổng quát + SP)
% ------------------------------------------------------------------------------
Tuy nhiên, việc ứng dụng các mô hình LLM thương mại vào nghiên cứu học thuật đang đối mặt với những thách thức nghiêm trọng: hiện tượng bịa đặt thông tin và trích dẫn giả mạo (Hallucination); thiếu cơ chế kiểm chứng trực quan tọa độ hộp bao (Bounding Box); rào cản trong biên soạn bài báo chuẩn quốc tế bằng LaTeX; và nguy cơ rò rỉ quyền sở hữu trí tuệ bản thảo khi gửi lên máy chủ công cộng.

Xuất phát từ nhu cầu thực tiễn của đội ngũ giảng viên, học giả và nhà nghiên cứu tại Trường Đại học VinUni về một giải pháp chuyên biệt, an toàn và tin cậy, đề tài \textit{``Xây dựng nền tảng AI hỗ trợ nghiên cứu và viết bài báo khoa học VinUni dựa trên Agentic RAG, truy xuất tri thức học thuật và kiểm chứng trích dẫn''} được lựa chọn. Đề tài kết hợp Kiến trúc Lục giác, tìm kiếm lai pgvector kết hợp BM25, đồ thị tri thức GraphRAG, vòng phản tư ReAct tự sửa lỗi LaTeX trong sandbox E2B và bảo mật Zero Data Retention kết hợp BYOK. Đây là giải pháp hoàn chỉnh có giá trị khoa học cao và tính ứng dụng thực tiễn vượt trội.

\vspace{0.25cm}
\noindent \textbf{3. Mục tiêu}
\vspace{0.08cm}

\noindent \textbf{Mục tiêu tổng quát:}
\begin{itemize}[leftmargin=*, noitemsep, topsep=1pt, parsep=0pt]
    \item Xây dựng thành công ứng dụng web (VeriScholar) hoàn chỉnh, bảo đảm độ tin cậy và liêm chính học thuật tuyệt đối cho nghiên cứu khoa học.
    \item Hỗ trợ toàn trình quy trình công bố: từ tra cứu tài liệu học thuật uy tín, đọc hiểu sâu có kiểm chứng tọa độ trực quan, tổng quan y văn qua đồ thị tri thức, đến soạn thảo bài báo LaTeX tự phục hồi lỗi biên dịch.
\end{itemize}

\vspace{0.15cm}
\noindent \textbf{Mục tiêu cụ thể:}

\noindent \textbf{-- Về sản phẩm:}
\begin{itemize}[leftmargin=*, noitemsep, topsep=1pt, parsep=0pt]
    \item Phân hệ Khám phá tri thức \& Thu thập đa nguồn: Tích hợp tìm kiếm học thuật hai tầng (Exa AI, Semantic Scholar, CrossRef, arXiv) với bộ lọc ngữ nghĩa nghiêm ngặt.
    \item Phân hệ Đọc sâu \& Bảo chứng trích dẫn tọa độ: Trình đọc đa tài liệu PDF Canvas kết hợp tìm kiếm lai (pgvector HNSW + BM25, dung hợp RRF) và trực quan hóa khung viền tọa độ (Bounding Box).
    \item Phân hệ Đồ thị tri thức học thuật (GraphRAG): Mô hình hóa 4 loại thực thể khoa học (Khái niệm, Phương pháp, Tiêu chuẩn đối chuẩn, Tác giả) và phát hiện cụm tri thức cộng đồng.
    \item Phân hệ Biên soạn trực tuyến \& Tác nhân AI sửa lỗi: Môi trường soạn thảo bài báo LaTeX chuẩn IEEE/ACM tích hợp sandbox máy ảo cô lập E2B và tác nhân ReAct tự động sửa lỗi biên dịch.
    \item Phân hệ Bảo mật chủ quyền \& Liêm chính nghiên cứu: Cơ chế mã hóa khóa riêng BYOK phía máy khách và nguyên tắc không lưu trữ dữ liệu (Zero Data Retention) bảo vệ tuyệt đối ý tưởng học thuật.
\end{itemize}

\newpage
\thispagestyle{adminform}

% ------------------------------------------------------------------------------
% TRANG 3 / 5 (Trang viii): Mục tiêu (Kỹ năng), 4. Nhiệm vụ & 5. Kế hoạch (Tuần 1)
% ------------------------------------------------------------------------------
\noindent \textbf{-- Về kiến thức và kỹ năng:}
\begin{itemize}[leftmargin=*, noitemsep, topsep=1pt, parsep=0pt]
    \item Vận dụng và làm sâu sắc kiến thức chuyên ngành: Kỹ nghệ phần mềm, CSDL nâng cao, Kiến trúc hệ thống, Học máy và Xử lý ngôn ngữ tự nhiên.
    \item Nắm vững và thực hành quy trình phát triển phần mềm theo phương pháp TDD, bảo đảm 100\% các cổng xác minh chất lượng trước khi bàn giao.
    \item Thành thạo các công nghệ hiện đại: FastAPI, SQLAlchemy Async, PostgreSQL/\allowbreak pgvector, React 19, TypeScript, Langfuse, LiteLLM, E2B MicroVM.
    \item Rèn luyện năng lực nghiên cứu độc lập, phân tích tài liệu quốc tế và tư duy giải quyết bài toán liêm chính học thuật.
\end{itemize}

\vspace{0.25cm}
\noindent \textbf{4. Các nhiệm vụ cụ thể}
\vspace{0.08cm}

Để hoàn thành mục tiêu đề ra, đồ án thực hiện các nhiệm vụ cụ thể sau:
\begin{enumerate}[leftmargin=*, noitemsep, topsep=1pt, parsep=0pt]
    \item Khảo sát bài toán nghiên cứu tại VinUni, thu thập tài liệu khoa học và lập đề cương đồ án.
    \item Phân tích và đặc tả các yêu cầu chức năng, phi chức năng của hệ thống (BYOK, IoU $\ge 0{,}85$).
    \item Thiết kế kiến trúc Lục giác, mô hình dữ liệu quan hệ và cơ sở dữ liệu véc-tơ pgvector.
    \item Thiết kế trải nghiệm người dùng (UI/UX) và các biểu đồ tuần tự nghiệp vụ cốt lõi.
    \item Xây dựng các module Backend cốt lõi: Xác thực, phân quyền, quản lý thư viện và lưu trữ CAS.
    \item Xây dựng phân hệ Hybrid RAG (pgvector Dense + BM25 Sparse) và kiểm chứng Bounding Box.
    \item Xây dựng phân hệ Đồ thị tri thức học thuật và tác nhân AI sửa lỗi LaTeX trong sandbox E2B.
    \item Xây dựng giao diện Frontend với Monaco Editor, trình xem PDF Canvas và streaming chat SSE.
    \item Thực hiện kiểm thử tự động TDD, đánh giá thực nghiệm và đóng gói sản phẩm bảo vệ.
\end{enumerate}

\vspace{0.25cm}
\noindent \textbf{5. Kế hoạch thực hiện}
\vspace{0.08cm}

\noindent
\begin{table}[h!]
\centering
\small
\setlength{\tabcolsep}{3.5pt}
\renewcommand{\arraystretch}{1.12}
\begin{tabular}{|>{\centering\arraybackslash}p{1.0cm}|>{\centering\arraybackslash}p{4.2cm}|>{\raggedright\arraybackslash}p{8.5cm}|}
\hline
\textbf{Tuần} & \textbf{Thời gian} & \textbf{Nội dung công việc} \\ \hline
1 & 24/08/2026 -- 30/08/2026 & Khảo sát bài toán nghiên cứu tại VinUni, thu thập tài liệu tham khảo liên quan, lập đề cương chi tiết và kế hoạch thực hiện đồ án. \\ \hline
\end{tabular}
\end{table}

\newpage
\thispagestyle{adminform}

% ------------------------------------------------------------------------------
% TRANG 4 / 5 (Trang ix): Kế hoạch thực hiện (Tuần 2 -- 9)
% ------------------------------------------------------------------------------
\noindent
\begin{table}[h!]
\centering
\small
\setlength{\tabcolsep}{3.5pt}
\renewcommand{\arraystretch}{1.08}
\begin{tabular}{|>{\centering\arraybackslash}p{1.0cm}|>{\centering\arraybackslash}p{4.2cm}|>{\raggedright\arraybackslash}p{8.5cm}|}
\hline
\textbf{Tuần} & \textbf{Thời gian} & \textbf{Nội dung công việc} \\ \hline
2 & 31/08/2026 -- 06/09/2026 & Phân tích yêu cầu nghiệp vụ; xác định các tác nhân (Actor), biểu đồ Use Case tổng thể và đặc tả chi tiết các Use Case chính. \\ \hline
3 & 07/09/2026 -- 13/09/2026 & Thiết kế kiến trúc Lục giác (Hexagonal Architecture), thiết kế CSDL quan hệ PostgreSQL và cơ sở dữ liệu véc-tơ pgvector. \\ \hline
4 & 14/09/2026 -- 20/09/2026 & Thiết kế giao diện người dùng khoa học tối giản (UI/UX); thiết kế biểu đồ tuần tự cho các luồng nghiệp vụ RAG và kiểm chứng Bounding Box. \\ \hline
5 & 21/09/2026 -- 27/09/2026 & Xây dựng Backend cốt lõi: Cấu hình môi trường, triển khai tầng Core nghiệp vụ, xác thực JWT, lưu trữ CAS và mã hóa BYOK. \\ \hline
6 & 28/09/2026 -- 04/10/2026 & Hiện thực hóa phân hệ Hybrid RAG: Tìm kiếm kết hợp HNSW + BM25, dung hợp RRF, tái xếp hạng Cross-Encoder và trích xuất Bounding Box PDF. \\ \hline
7 & 05/10/2026 -- 11/10/2026 & Hiện thực hóa phân hệ Đồ thị tri thức (GraphRAG) và Tác nhân AI ReAct tự động phục hồi lỗi biên dịch LaTeX trong sandbox E2B. \\ \hline
8 & 12/10/2026 -- 18/10/2026 & Xây dựng giao diện Frontend (React 19): Monaco Editor, trình hiển thị PDF Canvas thời gian thực, bảng dẫn chứng Bbox và luồng SSE. \\ \hline
9 & 19/10/2026 -- 25/10/2026 & Tích hợp toàn diện hệ thống Frontend -- Backend; thực hiện kiểm thử tự động TDD (Unit Test, Boundary Test, E2E); tối ưu hóa hiệu năng và bảo mật. \\ \hline
\end{tabular}
\end{table}

\newpage
\thispagestyle{adminform}

% ------------------------------------------------------------------------------
% TRANG 5 / 5 (Trang x): Kế hoạch thực hiện (Tuần 10), Nơi thực hiện & Chữ ký
% ------------------------------------------------------------------------------
\vspace*{0.2cm}
\noindent
\begin{table}[h!]
\centering
\small
\setlength{\tabcolsep}{3.5pt}
\renewcommand{\arraystretch}{1.15}
\begin{tabular}{|>{\centering\arraybackslash}p{1.0cm}|>{\centering\arraybackslash}p{4.2cm}|>{\raggedright\arraybackslash}p{8.5cm}|}
\hline
\textbf{Tuần} & \textbf{Thời gian} & \textbf{Nội dung công việc} \\ \hline
10 & 26/10/2026 -- 02/11/2026 & Đánh giá thực nghiệm trên các ấn phẩm khoa học thực tế; hoàn thiện báo cáo tốt nghiệp, slide thuyết trình, đóng gói sản phẩm và nộp báo cáo về Bộ môn Hệ thống Thông tin ngày 02/11/2026. \\ \hline
\end{tabular}
\end{table}

\vspace{0.4cm}
\noindent \textbf{6. Nơi thực hiện đồ án:} Trường Đại học Mỏ -- Địa chất

\vspace{1.8cm}
\hfill \textit{Hà Nội, ngày\dots\dots tháng\dots\dots năm 2026}
\vspace{0.5cm}

\begin{center}
\begin{tabular}{>{\centering\arraybackslash}p{7.0cm} >{\centering\arraybackslash}p{7.0cm}}
  \textbf{Giảng viên hướng dẫn} & \textbf{Sinh viên thực hiện} \\[0.3cm]
  \textit{(ký, họ tên)} & \textit{(ký, họ tên)} \\[2.8cm]
  \textbf{TS. Vũ Lan Phương} & \textbf{Đàm Mạnh Dũng} \\
\end{tabular}
\end{center}

} % kết thúc setstretch
